% --- Fonts ---
% Load Termes (Times) as the main text font
\fontfam[Termes] 

% --- Legacy Compatibility (The Shim) ---
\def\twelvesl{\it}  % Old 'slanted' -> Italic
\def\twelveit{\it}  % Old 'italic' -> Italic
\def\twelvebf{\bf}  % Old 'bold' -> Bold
\def\twelverm{\rm}  % Old 'roman' -> Normal / Reset to Termes

% Fix for Sans Serif (ss):
% Switch to Heros (Helvetica) whenever 'ss' is requested.
\def\twelvess{\fontfam[Heros]} 
\def\seventeenss{\typosize[17/21]\fontfam[Heros]}

% Fix for Typewriter (tt):
% Switch to Cursor (Courier) whenever 'tt' is requested.
\def\twelvett{\fontfam[Cursor]} 

% --- Legacy Compatibility (The Shim) ---
...
\def\tentt{\fontfam[Cursor]} 

% FIX FOR VERTICAL SPACING ALERTS:
\raggedbottom

% FIX FOR \ttfam CRASH:
% Define \ttfam as 0. This stops the crash.
% (Note: Numbers in math mode might look like Roman text instead of Typewriter, 
% but the document will compile).
\chardef\ttfam=0

% Set base font size to 10pt with 12pt line spacing
\typosize[10/12]

% --- Language (UK English) ---
\enlang \uselanguage{ukenglish}

% --- Page Dimensions & Margins ---
\margins/2 (129,198) (15,15,18.5,18.5)mm

% --- Paragraph Settings ---
\parskip=0pt
\parindent=12pt
\tolerance=1414
\hbadness=1414
\emergencystretch=5pt

% --- Fonts (The Faber Voice) ---
% Faber uses Garamond/Bembo. Termes (Times) is okay, but Garamond is the "Disguise".
% Un-comment the next line if you installed EB Garamond, otherwise stick to Termes.
\fontfam[EBGaramond] 
%\fontfam[Termes] 

% Define the Display Font (Albertus) for Titles/Chapters
% \font\Albertus=[Ferry-Regular] at 18pt
\font\Albertus=[texgyreheros-bold] at 18pt

% --- Legacy Compatibility Shims (KEPT AS IS) ---
\def\twelvesl{\it}  
\def\twelveit{\it}  
\def\twelvebf{\bf}  
\def\twelverm{\rm}  
\def\twelvess{\fontfam[Heros]} 
\def\seventeenss{\typosize[17/21]\fontfam[Heros]}
\def\twelvett{\fontfam[Cursor]} 
\def\tentt{\fontfam[Cursor]} 
\chardef\ttfam=0 
\raggedbottom

% --- Front Matter Layout Macros ---

% Semantic vertical spacing (strictly metric)
\def\frontmattergap{\vskip 10mm}
\def\titlegap{\vskip 25mm}
\def\publishergap{\vfill}
\def\copyrightgap{\vskip 4mm}

% A clean macro for intentionally blank pages
\def\blankpage{\null\vfill\eject}

% An environment that automatically centers all lines within it
\def\startcentered{%
   \begingroup
   \leftskip=0pt plus 1fill
   \rightskip=0pt plus 1fill
   \parindent=0pt
   \obeylines
}
\def\stopcentered{%
   \par
   \endgroup
   \vfill\eject
}

% --- Typography Settings ---
\typosize[10/13] % Slightly more air (13pt leading) feels more "bookish"
\parskip=0pt
\parindent=1.5em % Classic fiction indent
\tolerance=2500
\emergencystretch=20pt
\hfuzz=1pt % Prevents TeX from complaining about tiny overlaps
\hbadness=1414

% --- Language ---
\enlang \uselanguage{ukenglish}

% --- Page Dimensions (The "Golden Canon") ---
% Your previous margins (15mm) were a bit tight. 
% Faber books use deep bottom margins and wider outer edges.
\margins/2 (129,198) (18,18,22,25)mm

% --- Headers & Footers (The Faber Look) ---
\def\booktitle{ELEVEN SHORT STORIES ABOUT LOVE} % <--- VERIFY THIS
\def\authorname{PAUL M.~CRAY}

\nopagenumbers % Turn off default numbers so we can build custom ones

\headline={%
   \ifodd\pageno 
      \hfil {\caps \booktitle} \hfil % Recto: Title
   \else
      \hfil {\caps \authorname} \hfil % Verso: Author
   \fi
}

\footline={\hfil \the\pageno \hfil} % Page number at bottom center

% --- Chapter Styling Hook ---
% IMPORTANT: How do your tm1.tex files start?
% If they start with \chap or \chapter, we can style them here.
% If they use a custom command like \mychapter, tell me!

% --- First Paragraph Indent Suppressor ---
\def\suppressindent{%
   \par\nobreak
   \everypar={\setbox0=\lastbox \everypar={}}%
}

\def\printchap#1{%
   \vfil\break
   \begingroup
      \hyphenpenalty=10000 
      \null\vskip 30mm 
      \centerline{\Albertus\typosize[18/22] #1} 
      \vskip 15mm
      \nobreak
   \endgroup
   \noindent
}

% --- Chapter Styling Hook ---
\def\printchap#1#2{%
   \vfil\break
   \begingroup
      \hyphenpenalty=10000 
      \null\vskip 30mm 
      \centerline{\Albertus #1} 
      \vskip 5mm
      \centerline{\twelvett #2} 
      \vskip 15mm
      \nobreak
   \endgroup
   \suppressindent
}

% --- Draft Mode Switch ---
\newif\ifdraft
\drafttrue % Change to \draftfalse when you want to print the final book

% --- Chapter Metadata ---
\def\chaptermeta#1{%
   \ifdraft
      \vskip 2mm
      \centerline{\twelvess #1}
      \vskip 4mm
   \fi
   \suppressindent
}

% --- Telegram Environment ---
\def\telegramgap{\vskip 4mm} % Standard metric gap for inserts

\def\starttelegram{%
   \par\telegramgap
   \begingroup
   \twelvett % Or whatever typewriter font you defined
   \obeylines
   \leftskip=0pt plus 1fill
   \rightskip=0pt plus 1fill
   \parindent=0pt
}
\def\stoptelegram{%
   \par
   \endgroup
   \telegramgap
   \noindent % Ensures the text following the telegram isn't indented
}

% --- Email & Text Message Environment ---
\def\messagegap{\vskip 4mm}

\def\startmessage{%
   \par\messagegap
   \begingroup
   \twelvett % Switches to your Cursor (Courier) font
   \leftskip=12mm % Indents the block from the left margin
   \rightskip=12mm plus 1fill % Indents from the right AND makes it ragged-right
   \parindent=0pt % Removes the first-line indent for message blocks
   \parskip=2mm % Adds a little breathing room between paragraphs within the same message
}
\def\stopmessage{%
   \par
   \endgroup
   \messagegap
   \noindent % Ensures the narrative text following the message isn't indented
}

% --- Postcard Environment ---
\def\postcardgap{\vskip 6mm}

\def\startpostcard{%
   \par\postcardgap
   \begingroup
   \fontfam[Chorus]\it % TeX Gyre Chorus (requires \it to activate the Chancery script)
   \typosize[13/16] % Make it slightly larger as handwriting needs more breathing room
   \leftskip=15mm % Indent further than the emails to look like a physical card centered on the page
   \rightskip=15mm plus 1fill % Ragged right for natural handwriting flow
   \parindent=0pt
   \parskip=3mm
}
\def\stoppostcard{%
   \par
   \endgroup
   \postcardgap
   \noindent
}

% --- Receipt Environment ---
\def\receiptgap{\vskip 6mm}

\def\startreceipt{%
   \par\receiptgap
   \moveright 4mm \vbox\bgroup % Shifts the 85mm box slightly off the left margin
   \hsize=85mm 
   \typosize[9/11] 
   \fontfam[Cursor] 
   \parindent=0pt
}

\def\stopreceipt{%
   \egroup 
   \receiptgap
   \noindent
}

% Helper macro: Description gets 45mm, leaving a full 40mm for the price
\def\receiptitem#1#2{%
   \line{\vtop{\hsize=45mm \raggedright #1}\hfil #2}%
   \vskip 1mm
}

\def\receiptdash{\vskip 2mm \line{\xleaders\hbox{-}\hfil} \vskip 2mm}

% --- Document Content ---
\begingroup
  % Front matter usually has no headers/numbers
  \headline={} 
  \footline={}
  \input 11SSAL_front_matter.tex
\endgroup

% --- Scene Break (Whitespace Only) ---
\def\scenebreak{%
   \par % End the current paragraph
   \penalty-50 % Tells TeX this is a good place to break a page if it needs to
   \vskip 10mm plus 3mm minus 2mm % Flexible metric spacing (avoids bad vertical stretching)
   \suppressindent % Automatically removes the indent from the next paragraph
}

% --- Major Scene Break (With Asterisks) ---
\def\majorbreak{%
   \par
   \penalty-50
   \vskip 6mm plus 2mm minus 1mm
   \centerline{* \quad * \quad *} % Three cleanly spaced asterisks
   \vskip 6mm plus 2mm minus 1mm
   \suppressindent
}

% --- Business Card Environment ---
\def\cardgap{\vskip 8mm}

\def\startcard{%
   \par\cardgap
   \begingroup
   \fontfam[Heros] % Sans-serif (matches your \twelvess definition)
   \typosize[9/12] % Dropping it a point below body text gives it a neat, contained look
   \obeylines % Respects your physical line breaks
   \leftskip=0pt plus 1fill % Pushes from the left
   \rightskip=0pt plus 1fill % Pushes from the right (creating an auto-center)
   \parindent=0pt
}

\def\stopcard{%
   \par
   \endgroup
   \cardgap
   \noindent
}

% --- UI Button Font & Macro ---
% Loads Heros directly as a single font (not a whole family context)
\font\buttonfont=[texgyreheros-regular] at 9pt 

% The macro to use inline
\def\button#1{{\buttonfont #1}}

\pageno=1

\input 11SSAL_chapter1.tex
\input 11SSAL_chapter2.tex
\input 11SSAL_chapter3.tex
\input 11SSAL_chapter4.tex
\input 11SSAL_chapter5.tex
\input 11SSAL_chapter6.tex
\input 11SSAL_chapter7.tex
\input 11SSAL_chapter8.tex
\input 11SSAL_chapter9.tex
\input 11SSAL_chapter10.tex
\input 11SSAL_chapter11.tex

\bye
